% ============================================================
%  INTRODUCTION GÉNÉRALE
% ============================================================
\chapter*{Introduction générale}
\addcontentsline{toc}{chapter}{Introduction générale}
\markboth{Introduction générale}{Introduction générale}

% ======= COLLER ICI LE TEXTE DE L'INTRODUCTION GÉNÉRALE =======
% (copier depuis Canva)

\noindent [Texte de l'introduction générale]

\vspace{0.5cm}

\noindent Ce mémoire est organisé en quatre chapitres :

\begin{itemize}
  \item \textbf{Chapitre 1 — Cadre général :} présente le contexte, la problématique, la solution proposée et la méthodologie Scrum adoptée.
  \item \textbf{Chapitre 2 — Conception 3D et identité visuelle :} décrit la modélisation 3D des produits sous Blender et les choix graphiques de la marque.
  \item \textbf{Chapitre 3 — Analyse et conception :} expose l'analyse des besoins, les diagrammes UML et l'architecture technique de la solution.
  \item \textbf{Chapitre 4 — Réalisation :} présente les interfaces développées, les fonctionnalités implémentées et les résultats obtenus.
\end{itemize}
